\documentclass[12pt,a4paper]{article}
\usepackage{cmap}
\usepackage[utf8]{inputenc}
\usepackage[T2A]{fontenc}
\usepackage[russian]{babel}
\usepackage{amsmath,amssymb}
\usepackage{booktabs}
\newcommand{\csat}{c_o^{\mathrm{sat}}}
\newcommand{\bsat}{\beta_o^{\mathrm{sat}}}
\begin{document}

% Проверка сборки вставок плана BUG-019 в numerical_solver/doc/problem_statement.tex

\begin{tabular}{lll}
$\csat$       & полная сжимаемость нефти вдоль кривой насыщения & Па$^{-1}$ \\
$\bsat$       & полный коэфф.\ расширения $B_o$ вдоль кривой насыщения & К$^{-1}$ \\
$c_{ot}$      & сжимаемость насыщенной нефти вместе с~выделяющимся газом~\eqref{eq:c_ot} & Па$^{-1}$ \\
\end{tabular}

\begin{equation}\label{eq:co_cg}
  c_o = 1, \qquad c_g = 1
\end{equation}
\begin{equation}\label{eq:Cp}
  C_p = 1
\end{equation}

\textbf{Условие согласованности PVT.} При снижении давления насыщенная
нефть сжимается (уменьшается~$B_o$) и~выделяет газ объёмом
$B_g\,\Delta R_s$. Сжимаемость нефти вместе с~выделяющимся газом
\begin{equation}\label{eq:c_ot}
  \boxed{
    c_{ot} = \csat + \frac{B_g}{B_o}\,\frac{\partial R_s^{\text{eq}}}{\partial p}
    = \frac{1}{B_o}\left(B_g\,\frac{\partial R_s^{\text{eq}}}{\partial p}
      - \frac{d B_o^{\text{sat}}}{d p}\right) \ge 0,
  }
\end{equation}
где $B_o^{\text{sat}}(p) = B_o\bigl(p, R_s^{\text{eq}}(p)\bigr)$, должна быть
неотрицательной. Через неё коэффициент при~$\partial p/\partial t$
в~уравнении давления~\eqref{eq:Cp} равен
\[
  C_p = \frac{\phi}{B_g}\,\bigl(c_r + S_o\,c_{ot} + S_g\,c_g\bigr),
\]
и~$C_p \ge 0$ при любой~$S_g$ тогда и~только тогда, когда
$c_{ot} \ge 0$ и~$c_g \ge 0$~\eqref{eq:co_cg} (при $c_r = 0$); для
недонасыщенной нефти то~же требуется от~$c_o$. При $c_{ot} < 0$ в~ячейках
с~малой~$S_g$ коэффициент $C_p < 0$: уравнение давления становится обратным
уравнением теплопроводности и~теряет корректность, неявный шаг вырождается
на~последовательности значений~$\Delta t$ (условия (A2)--(A3):
M.~Peszy\'nska, Transport in Porous Media, 2006, т.~63; K.\,H.~Coats,
SPE~49774, 1999). Эмпирические корреляции нарушают условие при экстраполяции
за~область своих данных: корреляция Стендинга с~$B_g$ идеального газа для
нефти с~$\gamma_g = 0.83$, $45^\circ$\,API при $T = 380$~К даёт
$c_{ot} < 0$ при $p > 34.5$~МПа, то~есть при $R_s > 351$~м$^3$ст/м$^3$ст
(данные Стендинга~--- до~$R_s = 254$~м$^3$ст/м$^3$ст); на~такой нефти
с~$R_{s,0} = 500$~м$^3$ст/м$^3$ст расчёт не~сходится или зависит от~шага
по~времени. Решатель проверяет $c_{ot}$, $c_g$ и~$c_o$ при сборке модели
в~30 точках диапазона давлений задачи и~выдаёт предупреждение; расчёт
продолжается (так поступает ECLIPSE).

С~учётом $S_o + S_g = 1$ коэффициент~\eqref{eq:Cp} равен
$C_p = \phi\,(c_r + S_o\,c_{ot} + S_g\,c_g)/B_g$, где $c_{ot}$~---
сжимаемость насыщенной нефти вместе с~выделяющимся газом~\eqref{eq:c_ot};
при $c_{ot} < 0$ уравнение~\eqref{eq:Cp} в~ячейках с~малой~$S_g$ некорректно.

\end{document}
